% Part: proof-theory
% Chapter: sequent-calculus
% Section: rules-G3c

\documentclass[../../../include/open-logic-section]{subfiles}

\begin{document}

\begin{table}
\begin{tabular}{@{}c@{}c@{}}
\multicolumn{2}{@{}c@{}}{સ્વયંસિદ્ધો}\\[2ex]
$!C, \Gamma \Sequent !C$ & $\lfalse, \Gamma \Sequent \Delta$
\\[2ex]
\multicolumn{2}{@{}c@{}}{તાર્કિક નિયમો}\\[2ex]
\Axiom$\lnot !A, \Gamma \fCenter !A $
\RightLabel{\LeftR{\lnot}}
\UnaryInf$\lnot !A, \Gamma \fCenter \Delta$
\DisplayProof
&
\Axiom$!A, \Gamma \fCenter$
\RightLabel{\RightR{\lnot}}
\UnaryInf$ \Gamma \fCenter \lnot !A $
\DisplayProof
\\[4ex]
\Axiom$ !A, !B, \Gamma \fCenter \Delta$
\RightLabel{\LeftR{\land}}
\UnaryInf$ !A \land !B, \Gamma \fCenter \Delta$
\DisplayProof
&
\Axiom$\Gamma \fCenter !A$
\Axiom$ \Gamma \fCenter !B$
\RightLabel{\RightR{\land}}
\BinaryInf$ \Gamma \fCenter !A \land !B $
\DisplayProof
\\[4ex]\Axiom$!A, \Gamma \fCenter \Delta$
\Axiom$!B, \Gamma \fCenter \Delta$
\RightLabel{\LeftR{\lor}}
\BinaryInf$!A \lor !B, \Gamma \fCenter \Delta$
\DisplayProof
&
\begin{tabular}[t]{@{}r@{}}
\Axiom$\Gamma \fCenter !A$
\RightLabel{\RightR{\lor}}
\UnaryInf$ \Gamma \fCenter !A \lor !B$
\DisplayProof
\\[3ex]
\Axiom$\Gamma \fCenter !B$
\RightLabel{\RightR{\lor}}
\UnaryInf$ \Gamma \fCenter !A \lor !B$
\DisplayProof\\[3ex]
\end{tabular}
\\[4ex]
\Axiom$!A \lif !B, \Gamma \fCenter !A$
\Axiom$!B, \Gamma \fCenter \Delta$
\RightLabel{\LeftR{\lif}}
\BinaryInf$!A \lif !B, \Gamma \fCenter \Delta$
\DisplayProof
&
\Axiom$!A, \Gamma \fCenter !B$
\RightLabel{\RightR{\lif}}
\UnaryInf$ \Gamma \fCenter !A \lif !B $
\DisplayProof
\\[4ex]
\Axiom$ !A(t), \lforall[x][!A(x)], \Gamma \fCenter \Delta$
\RightLabel{\LeftR{\lforall}}
\UnaryInf$ \lforall[x][!A(x)],\Gamma \fCenter \Delta$
\DisplayProof
&
\Axiom$ \Gamma \fCenter !A(a) $
\RightLabel{\RightR{\lforall}}
\UnaryInf$ \Gamma \fCenter \lforall[x][!A(x)]$
\DisplayProof
\\[4ex]
\Axiom$ !A(a), \Gamma \fCenter \Delta $
\RightLabel{\LeftR{\lexists}}
\UnaryInf$ \lexists[x][!A(x)], \Gamma \fCenter \Delta$
\DisplayProof
&
\Axiom$ \Gamma \fCenter !A(t) $
\RightLabel{\RightR{\lexists}}
\UnaryInf$ \Gamma \fCenter \lexists[x][!A(x)]$
\DisplayProof
\end{tabular}
\caption{અંતઃપ્રજ્ઞાવાદી તર્ક માટે \Log{G3i}ના નિયમો.
$\RightR{\forall}$ અને $\LeftR{\exists}$માં !!{constant}~$a$
ફલિત-સિક્વન્ટમાં આવવો નહીં. $\Gamma$ એ !!{formula}નો બહુગણ છે,
અને $\Delta$ ખાલી હોઈ શકે અથવા એક જ !!{formula} ધરાવી શકે.
સ્વયંસિદ્ધોમાં !!{formula}~$!C$ પરમાણ્વીય હોવું જોઈએ.
અલગ તંત્ર અંગે નોંધ: \Log{G1m} એ જમણા શિથિલીકરણ અને
$\lfalse, \Gamma \Sequent
\Delta$ સ્વરૂપનાં સ્વયંસિદ્ધો વિનાનું \Log{G1i} છે.}
\ollabel{tab:G3c}
\end{table}
\sourcecorrection{OLMD-334}{મૂળ સર્વપરિમાણકના ડાબા નિયમના આધારમાં મુખ્ય સૂત્ર અને સંદર્ભ વચ્ચેનો અલ્પવિરામ ઉમેર્યો છે.}
\sourcecorrection{OLMD-335}{મૂળ વિકલ્પના જમણા નિયમના આધારમાં બે સૂત્રો છે, જ્યારે આ જ કોષ્ટકની શરત ઉત્તરાંગમાં વધુમાં વધુ એક સૂત્રની છૂટ આપે છે. તેથી એક સૂત્રવાળા બે જુદા દાખલ-નિયમ દર્શાવ્યા છે: એકમાં પહેલું વિકલ્પક અને બીજામાં બીજું વિકલ્પક આધાર છે. આ સ્વરૂપ G1iના કોષ્ટક અને નિયમોના અર્થઘટન સાથે મેળ ખાય છે.}
\sourcecorrection{OLMD-336}{મૂળની છેલ્લી કૅપ્શન-નોંધ G1નાં તંત્રો અંગે છે; તેને આ G3i કોષ્ટકની વ્યાખ્યા ગણવી નહીં. G1iના અલગ કોષ્ટક અનુસાર G1mમાં જમણું શિથિલીકરણ પણ દૂર કરવું પડે છે. અહીં એ ચૂકી ગયેલી શરત ઉમેરેલી છે; ફક્ત અસત્યનાં સ્વયંસિદ્ધો દૂર કરવાથી મૂળ નકાર-નિયમો સાથે ઇચ્છિત ન્યૂનતમ તર્ક મળતું નથી.}

\end{document}
